\section{未来实验验证}
\label{sec:future}

本文不包含任何湿实验数据；本节仅说明如何用最小实验设计检验平台输出的候选假设。

\textbf{方向一：第 18 位碱基替换（候选 C1）.} 最小设计为：选择若干携带目标位点的 sgRNA，仅替换第 18 位碱基（C$\leftrightarrow$A），保持其余 22\,nt 完全不变，在同一实验体系中测定编辑效率。该设计直接检验"模型优先排序的核苷酸差异是否伴随编辑效率变化"，并可通过在多于一个细胞系中重复，检验平台提示的上下文依赖性（HCT116/HeLa 与 HEK293T 的方向差异）。读出指标建议使用与本文数据一致的归一化编辑效率（indel 比例或报告基因活性），并预先设定样本量与重复次数。

\textbf{方向二：PAM 邻近候选模式破坏（候选 C2/C3）.} 最小设计为：在包含 GAGG/GGGG 核心的 sgRNA 中替换该 4\,nt 核心（保持 GC 含量与长度不变以控制混杂），比较携带与不携带核心的序列的实测效率。由于该候选目前仅在 HeLa 中达到富集，实验应优先在 HeLa 中进行，并在其他细胞系中作为探索性重复。

\textbf{方向三：独立数据集与真正 LOCO 验证.} 本批次的留一细胞系分支退化，因此需要重新执行训练配置以完成真正的 LOCO 评估；此外，应在独立编辑数据集（不同细胞类型或不同编辑系统）上检验平台流程的可迁移性，包括候选模式与位置效应的可重复性。

\textbf{预期输出与判定标准.} 上述实验的预期输出为"候选假设被支持/不被支持"的二值判定与效应量估计；判定应基于预先登记的统计方案（分层比较与多重检验校正），并使用与本文一致的指标口径。无论结果方向如何，实验都将用于评估平台在"AI 候选优先级 $\rightarrow$ 实验验证"闭环中的实用价值。
